%% This is file `elsarticle-template-1-num.tex',
%%
%% Copyright 2009 Elsevier Ltd
%%
%% This file is part of the 'Elsarticle Bundle'.
%% ---------------------------------------------

\documentclass[a4paper]{article}

\usepackage{graphicx}
\usepackage{import}
\usepackage{xifthen}
\usepackage{pdfpages}
\usepackage{transparent}
\usepackage{cancel}
\usepackage{amssymb}
\usepackage{amsmath}
\usepackage{amsthm}
\usepackage{ mathrsfs }
\usepackage{ dsfont }
\usepackage{tikz-cd}
\usepackage[a4paper,top=3cm,bottom=2cm,left=3cm,right=3cm,marginparwidth=1.75cm]{geometry}
\usepackage{float}
\usepackage{lineno}
\usepackage{hyperref}
\usepackage{enumitem}

\newtheorem{name}{Printed output}
\newtheorem{mydef}{Definition}
\newtheoremstyle{break}%
  {}%
  {}%
  {}%
  {}%
  {\bfseries}%
  {.}%
  {\newline}%
  {}

\theoremstyle{plain}
\newtheorem{question}{Question}

\newtheorem{theorem}{Theorem}[section]
\newtheorem*{lemma*}{Lemma}
\newtheorem{lemma}[theorem]{Lemma}
\newtheorem{corollary}[theorem]{Corollary}
\newtheorem{proposition}[theorem]{Proposition}

\theoremstyle{definition}
\newtheorem{definition}[theorem]{Definition}
\newtheorem*{note*}{Note}
\newtheorem*{claim}{Claim}
\newtheorem*{remark}{Remark}
\newtheorem{example}{Example}
\newtheorem*{solution}{Solution}
\newtheorem{exercise}{Exercise}

\numberwithin{equation}{subsection}

\newcommand{\incfig}[2][1]{%
    \def\svgwidth{#1\columnwidth}%
    \import{./Figures/}{#2.pdf_tex}%
}
\newcommand{\hodge}{{*}}
\newcommand{\tang}{\operatorname{\textbf{\textup{t}}}}
\newcommand{\norm}{\operatorname{\textbf{\textup{n}}}}
\newcommand{\R}{\mathbb{R}}
\newcommand{\bP}{\mathbb{bP}}
\newcommand{\N}{\mathbb{N}}
\newcommand{\Q}{\mathbb{Q}}
\newcommand{\A}{\mathbb{A}}
\newcommand{\bS}{\mathbb{S}}
\newcommand{\Z}{\mathbb{Z}}
\newcommand{\F}{\mathbb{F}}
\newcommand{\B}{\mathbb{B}}
\newcommand{\T}{\mathbb{T}}
\newcommand{\K}{\mathbb{K}}
\newcommand{\C}{\mathbb{C}}
\newcommand{\D}{\mathbb{D}}

\usepackage{comment}
\newcommand{\bmv}{\bm{v}}
\newcommand{\bmu}{\bm{u}}
\newcommand{\bmw}{\bm{w}}
\newcommand{\bmx}{\bm{x}}

\newcommand{\cF}{\mathcal{F}}
\newcommand{\cA}{\mathcal{A}}
\newcommand{\cM}{\mathcal{M}}
\newcommand{\cU}{\mathcal{U}}
\newcommand{\cB}{\mathcal{B}}
\newcommand{\cD}{\mathcal{D}}
\newcommand{\cL}{\mathcal{L}}
\newcommand{\cS}{\mathcal{S}}
\newcommand{\cG}{\mathcal{G}}
\newcommand{\cT}{\mathcal{T}}
\newcommand{\cC}{\mathcal{C}}
\newcommand{\cH}{\mathcal{H}}
\newcommand{\cO}{\mathcal{O}}
\newcommand{\cJ}{\mathcal{J}}

\newcommand{\frp}{\mathfrak{p}}
\newcommand{\frq}{\mathfrak{q}}
\newcommand{\frm}{\mathfrak{m}}
\newcommand{\frP}{\mathfrak{P}}

\newcommand{\msF}{\mathscr{F}}
\newcommand{\msA}{\mathscr{A}}
\newcommand{\msM}{\mathscr{M}}
\newcommand{\msE}{\mathscr{E}}
\newcommand{\msU}{\mathscr{U}}
\newcommand{\msB}{\mathscr{B}}
\newcommand{\msD}{\mathscr{D}}
\newcommand{\msL}{\mathscr{L}}
\newcommand{\msS}{\mathscr{S}}
\newcommand{\msG}{\mathscr{G}}
\newcommand{\msT}{\mathscr{T}}
\newcommand{\msC}{\mathscr{C}}
\newcommand{\msH}{\mathscr{H}}
\newcommand{\msO}{\mathscr{O}}
\newcommand{\msJ}{\mathscr{J}}
\newcommand{\msP}{\mathscr{P}}

\newcommand{\vect}[1]{\underline{\textbf{#1}}}
\newcommand{\llangle}{\left\langle}
\newcommand{\rrangle}{\right\rangle}
\newcommand{\llbrace}{\left\lbrace}
\newcommand{\rrbrace}{\right\rbrace}
\newcommand{\setof}[1]{\left\lbrace #1 \right\rbrace}

\newcommand{\cFA}{\cF(\cA)}
\newcommand{\sigmaA}{\sigma(\cA)}
\newcommand{\del}[2]{\frac{\partial #1}{\partial #2}}

\newcommand{\innerproductb}[2]{\llangle\bm{#1},\bm{#2}\rrangle}
\newcommand{\innerproduct}[2]{\llangle{#1},{#2}\rrangle}

\newcommand{\cl}{\operatorname{cl}}
\newcommand{\coker}{\operatorname{coker}}
\newcommand{\re}{\operatorname{Re}}
\newcommand{\im}{\operatorname{Im}}
\newcommand{\dist}{\operatorname{dist}}
\newcommand{\Span}{\operatorname{Span}}
\newcommand{\Supp}{\operatorname{Supp}}
\newcommand{\id}{\operatorname{id}}
\newcommand{\sh}{\operatorname{sh}}
\newcommand{\var}{\operatorname{Var}}
\newcommand{\Proj}{\operatorname{Proj}}
\newcommand{\Ker}{\operatorname{Ker}}
\newcommand{\Mor}{\operatorname{Mor}}
\newcommand{\supp}{\operatorname{supp}}
\newcommand{\Ran}{\operatorname{Ran}}
\newcommand{\Res}{\operatorname{Res}}
\newcommand{\Mod}{\operatorname{Mod}}
\newcommand{\res}{\operatorname{res}}
\newcommand{\Frac}{\operatorname{Frac}}
\newcommand{\End}{\operatorname{End}}
\newcommand{\Sym}{\operatorname{Sym}}
\newcommand{\Spec}{\operatorname{Spec}}
\newcommand*{\sheafhom}{\mathcal{H}\kern -.5pt \it{ om}}
\newcommand*{\sheafspec}{\mathcal{S}\kern -.5pt \it{ pec}}
\newcommand*{\sheafproj}{{\mathcal{P}\kern -.5pt \it{ roj}}}
\newcommand{\Div}{\operatorname{div}}
\newcommand{\Aut}{\operatorname{Aut}}
\newcommand{\Hom}{\operatorname{Hom}}
\newcommand{\Qcoh}{\operatorname{Qcoh}}
\newcommand{\PGL}{\operatorname{PGL}}
\newcommand{\Op}{\operatorname{Op}}
\newcommand{\WF}{\operatorname{WF}}
\newcommand{\Char}{\operatorname{Char}}
\newcommand{\Ell}{\operatorname{Ell}}
\newcommand{\RE}{\operatorname{Re}}
\newcommand{\IM}{\operatorname{Im}}
\newcommand{\cosec}{\operatorname{cosec}}

\setlength{\parindent}{0pt}
\setlength{\parskip}{1em}

\begin{document}

\title{\textbf{Problem Sheet: Functions 1 Solution}}
\author{\textbf{Jeffrey Thompson}}
\maketitle

\begin{exercise}
Let $f$ be a function defined by
\[
    f(x) = \frac{3x+2}{7}.
\]
\begin{enumerate}
    \item Find $ff(x)$, and write your answer in the form
    \[
        ff(x) = \frac{ax+b}{49}
    \]
    where $a,b$ are numbers.
    \item Find $f^{-1}(x)$, and write your answer in the form
    \[
        f^{-1}(x) = \frac{ax+b}{3}
    \]
    where $a,b$ are numbers.
\end{enumerate}
\end{exercise}

\begin{solution}

\begin{enumerate}
    \item First write
    \[
        f(x) = \frac{3x+2}{7}.
    \]
    Then
    \[
        ff(x) = f(f(x)) = \frac{3f(x)+2}{7}.
    \]
    Substituting $f(x)$ in gives
    \[
        ff(x) = \frac{3\frac{3x+2}{7}+2}{7}.
    \]
    Multiply by $\frac{7}{7}$ to remove the inner denominator:
    \begin{align*}
        ff(x)
        &= \frac{7}{7}\times\frac{3\frac{3x+2}{7}+2}{7}\\
        &= \frac{3\times\cancel{7}\times\frac{3x+2}{\cancel{7}}+14}{49}\\
        &= \frac{9x+6+14}{49}\\
        &= \frac{9x+20}{49}.
    \end{align*}
    Hence
    \[
        ff(x) = \frac{9x+20}{49}
    \]
    and so $a=9$, $b=20$.

    \item Write
    \[
        y = \frac{3x+2}{7}.
    \]
    Rearranging for $x$:
    \begin{align*}
        7y &= 3x+2\\
        7y-2 &= 3x\\
        x &= \frac{7y-2}{3}.
    \end{align*}
    Therefore
    \[
        f^{-1}(x) = \frac{7x-2}{3}.
    \]
    Hence $a=7$, $b=-2$.
\end{enumerate}
\qed
\end{solution}

% ------------------------------------------------------------

\begin{exercise}
Let $f$ be a function defined by
\[
    f(x) = \frac{3x+2}{6x}, \qquad x \neq 0.
\]
\begin{enumerate}
    \item Find $ff(x)$, simplifying your answer into the form
    \[
        ff(x) = \frac{ax+b}{cx+d}
    \]
    where $a,b,c,d$ are numbers.
    \item State where $ff(x)$ is defined.
    \item Find $f^{-1}(x)$ in the form
    \[
        f^{-1}(x) = \frac{1}{ax+b}
    \]
    where $a,b$ are numbers.
    \item State where $f^{-1}(x)$ is defined.
\end{enumerate}
\end{exercise}

\begin{solution}

\begin{enumerate}
    \item First write
    \[
        f(x) = \frac{3x+2}{6x}, \qquad x \neq 0.
    \]
    Then
    \[
        ff(x) = f(f(x)) = \frac{3f(x)+2}{6f(x)}.
    \]
    Substituting gives
    \[
        ff(x) = \frac{3\frac{3x+2}{6x}+2}{6\frac{3x+2}{6x}}.
    \]
    Multiply by $\frac{6x}{6x}$:
    \begin{align*}
        ff(x)
        &= \frac{6x}{6x}\times \frac{3\frac{3x+2}{6x}+2}{6\frac{3x+2}{6x}}\\
        &= \frac{3\times\cancel{6x}\times\frac{3x+2}{\cancel{6x}}+12x}{6\times\cancel{6x}\times\frac{3x+2}{\cancel{6x}}}\\
        &= \frac{9x+6+12x}{18x+12}\\
        &= \frac{21x+6}{18x+12}\\
        &= \frac{7x+2}{6x+4}.
    \end{align*}
    Hence
    \[
        ff(x) = \frac{7x+2}{6x+4}
    \]
    so $a=7$, $b=2$, $c=6$, $d=4$.

    \item We need
    \[
        6x+4 \neq 0
    \]
    so
    \[
        x \neq -\frac{2}{3}.
    \]
    But we also need $x \neq 0$ from the original function. There is one more thing to check: since $ff(x)=f(f(x))$, we also need
    \[
        f(x) \neq 0.
    \]
    That would mean
    \[
        \frac{3x+2}{6x} = 0,
    \]
    but this has no solution, since a fraction is only $0$ when its numerator is $0$, and $3x+2=0$ would give $x=-\frac{2}{3}$, which we have already excluded from the denominator of $ff(x)$.
    Therefore $ff(x)$ is defined whenever
    \[
        x \neq 0, \qquad x \neq -\frac{2}{3}.
    \]

    \item Write
    \[
        y = \frac{3x+2}{6x}.
    \]
    Rearranging for $x$:
    \begin{align*}
        6xy &= 3x+2\\
        6xy-3x &= 2\\
        x(6y-3) &= 2\\
        x &= \frac{2}{6y-3}\\
        &= \frac{1}{3y-\frac{3}{2}}.
    \end{align*}
    Therefore
    \[
        f^{-1}(x) = \frac{1}{3x-\frac{3}{2}}.
    \]
    Hence $a=3$, $b=-\frac{3}{2}$.

    \item We need
    \[
        3x-\frac{3}{2} \neq 0
    \]
    and therefore
    \[
        x \neq \frac{1}{2}.
    \]
\end{enumerate}
\qed
\end{solution}

% ------------------------------------------------------------

\begin{exercise}
Let $f$ and $g$ be functions defined by
\[
    f(x) = \sin(x) + \cos(2x)
\]
\[
    g(x) = \frac{x-1}{2}.
\]
\begin{enumerate}
    \item Find $fg(x)$, simplifying your answer.
    \item Let $h(x) = \frac{3x+2}{7}$ and $k(x) = \frac{7x-2}{3}$. By computing $hk(x)$, verify that $k$ is the inverse of $h$.
\end{enumerate}
\end{exercise}

\begin{solution}

\begin{enumerate}
    \item First write
    \[
        f(x) = \sin(x) + \cos(2x), \qquad g(x) = \frac{x-1}{2}.
    \]
    Then
    \[
        fg(x) = f(g(x)) = \sin(g(x)) + \cos(2g(x)).
    \]
    Substituting $g(x)$ gives
    \begin{align*}
        fg(x)
        &= \sin\left(\frac{x-1}{2}\right) + \cos\left(2\times\frac{x-1}{2}\right)\\
        &= \sin\left(\frac{x-1}{2}\right) + \cos(x-1).
    \end{align*}

    \item First write
    \[
        h(x) = \frac{3x+2}{7}, \qquad k(x) = \frac{7x-2}{3}.
    \]
    Then
    \begin{align*}
        h(k(x))
        &= \frac{3k(x)+2}{7}\\
        &= \frac{3\frac{7x-2}{3}+2}{7}\\
        &= \frac{\cancel{3}\frac{7x-2}{\cancel{3}}+2}{7}\\
        &= \frac{7x-2+2}{7}\\
        &= \frac{7x}{7}\\
        &= x.
    \end{align*}
    This already shows the key inverse relation. For completeness, we can also check the other composition:
    \begin{align*}
        k(h(x))
        &= \frac{7h(x)-2}{3}\\
        &= \frac{7\left(\frac{3x+2}{7}\right)-2}{3}\\
        &= \frac{3x+2-2}{3}\\
        &= \frac{3x}{3}\\
        &= x.
    \end{align*}
    Therefore both $h(k(x))=x$ and $k(h(x))=x$, which verifies that $k$ is the inverse of $h$.
\end{enumerate}
\qed
\end{solution}

\end{document}
